\documentclass{article}

\usepackage{fullpage}
\usepackage{amsmath}
\usepackage{amssymb}
\usepackage{xcolor}

\newcounter{placeholderCount}
\newcommand{\placeholder}[1]{{
    \stepcounter{placeholderCount}
    \color{red} [#1]
}}

\newcounter{solutionCount}
\newcommand{\solution}[1]{
    \stepcounter{solutionCount}
    \paragraph{Solution \thesolutionCount.}
    {
        #1
    }
}

\title{
    Deep Learning (Spring 2026) \\
    Solution for Assignment 3
}
\author{
    Pan Changxun
    \\
    2024011323
}

\begin{document}

\maketitle

\solution{
    False. FID (Fréchet Inception Distance) measures the distance between the feature distributions of real and generated images. A smaller distance means higher similarity to real images. Therefore, a better generator achieves a lower FID score, not a higher one.
}

\solution{
    False. Training the discriminator to absolute optimality before updating the generator leads to the vanishing gradient problem. A perfect discriminator will reject fake images with near-absolute certainty, providing essentially zero gradients for the generator to learn from. In practice, they are trained using alternating updates.
}

\solution{
    \begin{enumerate}
        \item Since $JSD(p||q) = \frac{1}{2} \left( KL(p||m) + KL(q||m) \right)$, where $m = \frac{1}{2}(p+q)$, we express it using log probabilities:
        \[
        JSD(p||q) = \frac{1}{2} \sum_x p(x) \log\frac{p(x)}{m(x)} + \frac{1}{2} \sum_x q(x) \log\frac{q(x)}{m(x)}.
        \]
        Also, $m(x) = \frac{1}{2}(p(x) + q(x))$. Using Shannon Entropy $H(p) = -\sum_x p(x) \log p(x)$, we can rewrite:
        \begin{align*}
            JSD(p||q) &= \frac{1}{2} \sum_x p(x) \log\frac{p(x)}{\frac{1}{2}(p(x)+q(x))} + \frac{1}{2} \sum_x q(x) \log\frac{q(x)}{\frac{1}{2}(p(x)+q(x))}\\
            &= \frac{1}{2} \sum_x p(x) \log p(x) + \frac{1}{2} \sum_x q(x) \log q(x) - \sum_x m(x) \log m(x)\\
            &= -\frac{1}{2} H(p) - \frac{1}{2} H(q) + H(m)\\
            &= H\!\left(\frac{p+q}{2}\right) - \frac{1}{2}(H(p) + H(q)).
        \end{align*}

        \item Now prove $0 \le JSD(p||q) \le \log 2$. For the lower bound, since KL divergence is non-negative (Gibbs' inequality), $JSD(p||q) \ge 0$. 
        
        For the upper bound, we use the fact that $KL(q||m) = \sum_x q(x) \log\frac{q(x)}{m(x)}$ and $m(x) = \frac{1}{2}(p(x)+q(x))$. The fraction $\frac{q(x)}{m(x)}=\frac{2q(x)}{q(x)+p(x)} \le 2$ implies $\log\frac{q(x)}{m(x)} \le \log 2$. Summing over the distribution yields $KL(q||m) \le \log 2$. Similarly, $KL(p||m) \le \log 2$. Thus:
        \[
        JSD(p||q) = \frac{1}{2} KL(p||m) + \frac{1}{2} KL(q||m) \le \frac{1}{2} \log 2 + \frac{1}{2} \log 2 = \log 2.
        \]
        When $q(x)\cdot p(x) = 0$ for all $x$ (disjoint supports), $JSD(p||q) = \log 2$. When $p = q$, $JSD(p||q) = 0$.

        \item To prove the triangle inequality $\sqrt{JSD(p_1||p_3)} \le \sqrt{JSD(p_1||p_2)} + \sqrt{JSD(p_2||p_3)}$, we demonstrate that $\sqrt{JSD}$ is a valid metric. 
        
        According to Endres and Schindelin (2003), the terms in the Jensen-Shannon divergence can be expressed using the following integral identity for any $p, q > 0$:
        \[
        p \log \frac{2p}{p+q} + q \log \frac{2q}{p+q} = \int_{0}^{\infty} \frac{(p - q)^2}{(p+t)(q+t)(p+q+2t)} dt
        \]
        Summing over all discrete states $x$, the JSD becomes:
        \[
        JSD(p||q) = \frac{1}{2} \sum_x \int_{0}^{\infty} \frac{(p(x) - q(x))^2}{(p(x)+t)(q(x)+t)(p(x)+q(x)+2t)} dt
        \]
        Because the integrand is scaled by a strictly positive weight $(p(x) - q(x))^2$, this integral representation proves that $JSD(p||q)$ is a negative definite kernel. Consequently, it corresponds to the squared Euclidean distance in some Hilbert space $\mathcal{H}$, such that $JSD(p||q) = \|\Phi(p) - \Phi(q)\|_{\mathcal{H}}^2$ for some feature map $\Phi$.
        
        Therefore, $\sqrt{JSD(p||q)}$ is the Euclidean distance $\|\Phi(p) - \Phi(q)\|_{\mathcal{H}}$, which inherently satisfies the Minkowski inequality (triangle inequality):
        \[
        \|\Phi(p_1) - \Phi(p_3)\|_{\mathcal{H}} \le \|\Phi(p_1) - \Phi(p_2)\|_{\mathcal{H}} + \|\Phi(p_2) - \Phi(p_3)\|_{\mathcal{H}}
        \]
        Substituting back the distance yields:
        \[
        \sqrt{JSD(p_1||p_3)} \le \sqrt{JSD(p_1||p_2)} + \sqrt{JSD(p_2||p_3)}.
        \]

    \end{enumerate}
}

\solution{
    \begin{enumerate}
        \item Using KR duality, we have:
        \[
        W(p, q) = \sup_{\|f\|_L \le 1} \left( \mathbb{E}_{x \sim p}[f(x)] - \mathbb{E}_{x \sim q}[f(x)] \right).
        \]
        Let \(f^*\) be the optimal function achieving the supremum. Then:
        \[
        W(p, q) = \mathbb{E}_{x \sim p}[f^*(x)] - \mathbb{E}_{x \sim q}[f^*(x)].
        \]
        consider another distribution \(r\). We can write:
        \[
        W(p, r) = \sup_{\|f\|_L \le 1} \left( \mathbb{E}_{x \sim p}[f(x)] - \mathbb{E}_{x \sim r}[f(x)] \right) \ge \mathbb{E}_{x \sim p}[f^*(x)] - \mathbb{E}_{x \sim r}[f^*(x)],
        \]
        and
        \[
        W(r, q) = \sup_{\|f\|_L \le 1} \left( \mathbb{E}_{x \sim r}[f(x)] - \mathbb{E}_{x \sim q}[f(x)] \right) \ge \mathbb{E}_{x \sim r}[f^*(x)] - \mathbb{E}_{x \sim q}[f^*(x)].
        \]
        Adding these two inequalities gives:
        \[
        W(p, r) + W(r, q) \ge \mathbb{E}_{x \sim p}[f^*(x)] - \mathbb{E}_{x \sim r}[f^*(x)] + \mathbb{E}_{x \sim r}[f^*(x)] - \mathbb{E}_{x \sim q}[f^*(x)] = W(p, q).
        \]
            Thus, we have shown that \(W(p, q) \le W(p, r) + W(r, q)\), proving the triangle inequality for the Wasserstein distance.
        \item Let's consider the Wasserstein distance between \(p_x\) and \(p_{x+\epsilon}\), choosing a simple \(\gamma^*\) that transports mass from \(x\) to \(x + \epsilon\):
        \[W(p_x, p_{x+\epsilon}) \le \sum_{x} \int_{\epsilon} p(x) p(\epsilon) \|x - (x + \epsilon)\|_2 d\epsilon = \int_{\epsilon} p(\epsilon) \|\epsilon\|_2 d\epsilon\]
        Namely the expected value of \(\|\epsilon\|_2\) under the distribution \(p(\epsilon)\). Using Cauthy Schwarz inequality, we have:
        \[
        \int_{\epsilon} p(\epsilon) \|\epsilon\|_2 d\epsilon \le \sqrt{\int_{\epsilon} p(\epsilon) d\epsilon} \cdot \sqrt{\int_{\epsilon} p(\epsilon) \|\epsilon\|_2^2 d\epsilon} = \sqrt{\mathbb{E}_{\epsilon \sim p(\epsilon)}[\|\epsilon\|_2^2]}.
        \]
        Bring them together, we have:
        \[
        W(p_x, p_{x+\epsilon}) \le \sqrt{\mathbb{E}_{\epsilon \sim p(\epsilon)}[\|\epsilon\|_2^2]}=V^{\frac{1}{2}}
        \]
        \item By the triangle inequality (part 1), we have:
        \[
        W(p_r, p_g) \le W(p_r, p_{r+\epsilon}) + W(p_{r+\epsilon}, p_{g+\epsilon}) + W(p_{g+\epsilon}, p_g).
        \]
        From part 2, $W(p_r, p_{r+\epsilon}) \le V^{1/2}$ and $W(p_g, p_{g+\epsilon}) \le V^{1/2}$. Thus:
        \[
        W(p_r, p_g) \le 2V^{1/2} + W(p_{r+\epsilon}, p_{g+\epsilon}).
        \]
        It remains to bound $W(p_{r+\epsilon}, p_{g+\epsilon})$. By KR duality:
        \[
        W(p_{r+\epsilon}, p_{g+\epsilon}) = \sup_{\|f\|_L \le 1} \left( \mathbb{E}_{p_{r+\epsilon}}[f] - \mathbb{E}_{p_{g+\epsilon}}[f] \right).
        \]
        Since both distributions have support on a ball of diameter $C$, any 1-Lipschitz function $f$ satisfies $0 \le f(x) \le C$ on the support (after shifting by a constant, which does not change the difference of expectations). Therefore:
        \begin{align*}
        \mathbb{E}_{p_{r+\epsilon}}[f] - \mathbb{E}_{p_{g+\epsilon}}[f] &= \sum_x f(x)\bigl(p_{r+\epsilon}(x) - p_{g+\epsilon}(x)\bigr) \\
        &\le C \sum_{x:\, p_{r+\epsilon}(x) > p_{g+\epsilon}(x)} \bigl(p_{r+\epsilon}(x) - p_{g+\epsilon}(x)\bigr) = C \cdot \mathrm{TV}(p_{r+\epsilon}, p_{g+\epsilon}).
        \end{align*}
        Next, we relate total variation to JSD. Let $m = \frac{1}{2}(p_{r+\epsilon} + p_{g+\epsilon})$. By Pinsker's inequality:
        \[
        \mathrm{TV}(p_{r+\epsilon},\, m) \le \sqrt{\tfrac{1}{2}\, KL(p_{r+\epsilon} \| m)}.
        \]
        Since $p_{r+\epsilon}(x) - m(x) = \frac{1}{2}(p_{r+\epsilon}(x) - p_{g+\epsilon}(x))$, we have $\mathrm{TV}(p_{r+\epsilon}, p_{g+\epsilon}) = 2\,\mathrm{TV}(p_{r+\epsilon}, m)$, so:
        \[
        \mathrm{TV}(p_{r+\epsilon}, p_{g+\epsilon}) \le 2\sqrt{\tfrac{1}{2}\, KL(p_{r+\epsilon}\|m)} \le 2\sqrt{JSD(p_{r+\epsilon}\|p_{g+\epsilon})},
        \]
        where the last step uses $KL(p_{r+\epsilon}\|m) \le 2\, JSD(p_{r+\epsilon}\|p_{g+\epsilon})$. Combining:
        \[
        W(p_r, p_g) \le 2V^{1/2} + 2C\sqrt{JSD(p_{r+\epsilon}\|p_{g+\epsilon})}.
        \]

        \item The result from part 3 suggests a practical trick: \textbf{add instance noise} to both real and generated images before feeding them to the discriminator. Specifically, we inject Gaussian noise $\epsilon \sim \mathcal{N}(0, \sigma^2 I)$ to both inputs. In this way, the discriminator effectively minimizes $JSD(p_{r+\epsilon}\|p_{g+\epsilon})$, which by the inequality in part 3, also minimizes an upper bound on the Wasserstein distance $W(p_r, p_g)$.

        Moreover, adding noise ensures that $p_{r+\epsilon}$ and $p_{g+\epsilon}$ have overlapping (full) support, which prevents the vanishing gradient problem that arises when $p_r$ and $p_g$ have disjoint supports.

        \textbf{Potential issue:} The noise introduces a bias term $2V^{1/2}$ in the upper bound. If $\sigma$ is too large, this bias dominates and the bound becomes loose, meaning the generator is not incentivized to match $p_r$ precisely. If $\sigma$ is too small, the smoothing effect is insufficient and JSD may still suffer from vanishing gradients. A common mitigation is to anneal (gradually decrease) the noise variance during training, but choosing the right schedule is non-trivial.
    \end{enumerate}
}


\solution{
    Used Claude (AI assistant by Anthropic) in the following capacities:
    \begin{itemize}
        \item \textbf{Solution ideation:} Brainstorming approaches and generating initial ideas for problem-solving strategies, particularly for structuring mathematical proofs and derivations.
        \item \textbf{English grammar and writing:} Polishing and refining the English writing throughout the solutions to improve clarity, coherence, and academic tone.
        \item \textbf{\LaTeX{} formatting:} Optimizing typesetting and layout of mathematical expressions, tables, and document structure for better readability and presentation.
        \item \textbf{Result verification:} Cross-checking the correctness of intermediate computations, final answers, and logical consistency of proofs.
    \end{itemize}
    All final solutions were reviewed and verified by the author.
}

\ifnum\value{placeholderCount}>0
    \vfill
    \begin{center}
        \color{red} \framebox{\textbf{WARNING: PLEASE REMOVE ALL PLACEHOLDERS BEFORE SUBMISSION}}
    \end{center}
\fi

\end{document}